\section{Stable frequency-assigned recovery in the nonlinear equations}
\label{sec:nonlinear-band}
The lower bounds at fixed snapshot times concern the information
record. Allowing each Fourier coefficient of a displacement field
to use a prescribed observation time changes that record. We now
give a nonlinear result valid in a regular-time neighborhood of
arbitrary smooth Gowdy data. The positive delay in the schedule is
independent of the spatial bandwidth.

Fix $a>0$, $s>1/2$ and a class of solutions on $[a,a+T]$ whose
$\|(P,Q)\|_{H^{s+1}}+\|(P_t,Q_t)\|_{H^s}$ are uniformly bounded.
Write $U=(P,Q)$ and write the coordinate field equations as
$U_{tt}-U_{\theta\theta}=\mathcal F(t,U,U_t,U_\theta)$,
where
\[
 \mathcal F=
 \begin{pmatrix}
 -P_t/t+e^{2P}(Q_t^2-Q_\theta^2)\\
 -Q_t/t-2(P_tQ_t-P_\theta Q_\theta)
 \end{pmatrix}.
\]
On this bounded class there is a Lipschitz constant $L$ such that
for the difference $Y=U-\widehat U$ and source difference $F$,
\[
 \|F(t)\|_{H^s}\le L E_s(Y(t)),\qquad
 E_s(Y)^2=\|Y_t\|_{H^s}^2+\|Y_\theta\|_{H^s}^2+\|Y\|_{H^s}^2.
\]
This follows from the algebra property of $H^s(\mathbb T)$ and
the Sobolev composition estimate for $e^{2P}$ on bounded sets.
For common initial $U(a)$, $E_s(Y(a))=\|Y_t(a)\|_{H^s}$.
The energy identity and Gronwall give
$E_s(Y(a+r))\le e^{(L+1)r}E_s(Y(a))$.
In particular let
\[
 \varepsilon_T = \frac{L}{L+1}\big(e^{(L+1)T}-1\big).
\]

For each recorded Fourier frequency $k$, use a delay $d_k\in(0,T]$
with $d_{-k}=d_k$ and set
\[
 w_k=\max\{|k|,1/d_k\},\qquad
 g_k=w_k S_k(d_k),\qquad
 S_k(r)=
 \begin{cases}
 \sin(|k|r)/|k|,& k\ne0,\\
 r,& k=0.
 \end{cases}
\]
We record $\widehat P(k,a+d_k)$, with the same initial
$P(a),Q(a),Q_t(a)$ for all compared solutions.
For comparing records, multiply each measured difference by $w_k$
and use the $H^s$ norm of the resulting Fourier sequence.
For a decoder, subtract the known free displacement
$\cos(|k|d_k)\widehat P(k,a)$ before multiplying.
The analogous record for both displacement fields
allows both initial velocities to be unknown.

\begin{theorem}[Nonlinear sampling estimate]
\label{th:nonlinear-band}
Suppose $\inf_k |g_k|\ge\gamma>\varepsilon_T$ on the
recorded frequencies. For two Gowdy solutions in the
bounded class, with common initial position and common
initial $Q_t$, write $f=P_t(a)-\widehat P_t(a)$,
and let $\mathcal M$ be the weighted record. With all
frequencies observed,
\[
 (\gamma-\varepsilon_T)\|f\|_{H^s}
 \le \|\mathcal M(U)-\mathcal M(\widehat U)\|_{H^s}
 \le (1+\varepsilon_T)\|f\|_{H^s}.
\]
With only $|k|\le K$ observed, the lower estimate becomes
\[
 (\gamma-\varepsilon_T)\|\Pi_K f\|_{H^s}
 \le \|\mathcal M_K(U)-\mathcal M_K(\widehat U)\|_{H^s}
       +\varepsilon_T\|(1-\Pi_K)f\|_{H^s}.
\]
The same estimates hold for both missing initial velocity
fields when both displacements are read.
\end{theorem}

\begin{proof}
Duhamel gives the observed $P$-difference on mode $k$
as $S_k(d_k)f_k+\int_0^{d_k} S_k(d_k-r)F_{P,k}(a+r)\,dr$.
For $0\le r\le d_k$, $w_k|S_k(d_k-r)|\le1$.
Indeed the two possible factors in $w_k$ use either
$|\sin x|\le1$ or $|S_k(r)|\le r$.
For the weighted sequence of integrals, use Minkowski after
extending each integral by zero to the common interval $[0,T]$.
The entire remainder is bounded in $H^s$ by
$\int_0^T\|F(a+r)\|_{H^s}\,dr\le\varepsilon_T\|f\|_{H^s}$.
This argument requires no diagonal structure for $F$ and admits
arbitrary coupling between modes.
The diagonal free term has gain between $\gamma$ and one.
The triangle inequality gives the two bounds and the version
with the omitted velocity tail. Vector-valued data use
the same argument. The periodic Gowdy momentum constraint
restricts the set of compared Cauchy data and does not
change the estimate.
\end{proof}

\paragraph{An explicit schedule with positive minimum delay.}
Set $N=\lceil 2\pi/T\rceil$ and $d_0=T/4$.
Read frequencies $|k|<N$ at $d_0$.
Their free phases have magnitudes below $\pi/2$, giving
$|g_k|\ge\sin(1)$, including $g_0=1$.
Group the remaining frequencies as
$N2^j\le |k|<N2^{j+1}$ for $j\ge0$.
For each group use the baseline $d_0$ and the option
$d_0+\pi/(4N2^j)$. For each mode keep the reading at
whichever of the two times gives the larger
$|\sin(|k|d)|$. Their phase separation is in
$[\pi/4,\pi/2)$, so the chosen magnitude is at least
$\gamma=\sin(\pi/8)$. These modes have $w_k=|k|$.
All chosen times are in $[a+T/4,a+3T/8]$.
There are at most $1+\max(0,1+\lfloor\log_2(K/N)\rfloor)$
distinct times through bandwidth $K$, and one saved
coefficient per frequency in addition to the initial
data. Since $1/d_k\le4/T$, the observation norm is
equivalent to $H^{s+1}$ on the assembled displacement
sequence, with constants independent of $K$.
At fixed $a$ and a fixed bounded set of nearby Cauchy data,
local existence gives a common regular-time interval and a common
bound, so $T$ can be reduced until $\varepsilon_T<\gamma$.
This includes neighborhoods of non-geodesic-image
solutions and the rank-two examples above.

\paragraph{Stable decoding and truncated data.}
Regard the forward map as $D v+R(v)$, with $D$ the
diagonal multipliers $g_k$ and $R$ Lipschitz with
constant $\varepsilon_T$. The subtraction of the known
free displacement has been included. Let $B$ be a closed
convex ball of trial velocities intersected with the
affine momentum constraint, and suppose the same forward
bound holds for every point of $B$.
For $|k|\le K$ take trial velocities in a fixed affine
band-limited slice, with their high-frequency part held fixed.
For both the infinite and finite-dimensional cases,
\[
 v_{n+1}=\operatorname{proj}_{B}
 \big[v_n+D^{-1}(\mathcal M_{\mathrm{data}}-\mathcal M(v_n))\big]
\]
is a contraction with ratio at most $\varepsilon_T/\gamma$.
In the affine finite-band case the multiplier inverse in
this formula acts only on the measured modes.
The projection is in the relevant Hilbert $H^s$ norm.
Consistent data recover the true velocity in $B$.
For additive error $\eta$ in the weighted record, the
resulting velocity error is at most
$\|\eta\|_{H^s}/(\gamma-\varepsilon_T)$ whenever the
true velocity lies in $B$. The truncated comparison
bound gives the additional $H^s$ tail term when the
true velocities are not in the same band-limited slice.
For an a priori $H^{s+r}$ bound, this tail is of order
$K^{-r}$ for $r>0$.

\paragraph{Timing errors.}
For known actual observation delays the proof applies
as written to those delays. For unknown timing errors
$\eta_k$ with both the nominal and actual times in
$[a,a+T]$, set $\eta_*=\sup_k w_k|\eta_k|$ using the
nominal weights. The change in the nominally weighted
record has norm bounded by
\[
 \eta_*\left(\|P_t(a)\|_{H^s}+\|P_\theta(a)\|_{H^s}
           +\int_0^T\|\mathcal F_P(a+r)\|_{H^s}\,dr\right).
\]
To see this, bound each temporal difference by
$|\eta_k|\sup_t|\widehat P_t(k,t)|$ and apply the
cosine propagator formula to $P_t$.
The bound is independent of the number of slots.
At high frequencies it corresponds to a clock error
on the order of $1/K$ for a uniform comparison.

\paragraph{Why the number of delayed slots grows.}
Consider a class containing the constant wave map and
suppose a uniform energy-normalized inverse estimate
holds with one later reading per mode and $m$ possible
delays in a fixed interval $[d_{\min},T]$, where
$d_{\min}>0$. Differentiate along a polarized single-mode
Cauchy family at the constant solution. The linear scalar
transfer in energy normalization is
$\sqrt{a/(a+d)}\{\sin(kd)+O(1/k)\}$,
uniformly on the fixed interval. For $k\ge 1/d_{\min}$
the prescribed weight is exactly $k$.
Choose a fixed integer $\ell$ large enough for the error
to be smaller than a quarter of the required lower gain
on multiples of $\ell$. Partition $[0,1)^m$ into
$R^m$ boxes for a fixed sufficiently large integer $R$.
Among $R^m+1$ multiples of the vector
$(\ell d_i/\pi)_{i=1}^m$ modulo one, two
belong to the same box. Their index difference gives
$k\in[\ell,\ell R^m]$ with $|\sin(kd_i)|\le\pi/R$
at every slot. This contradicts the required gain if
$\ell R^m\le K$. Thus the order $\log K$ is necessary
for uniform modewise energy conditioning in this class
when one reading per frequency is saved. The same
argument applies if the maximum of all available
readings per mode is used.

\paragraph{Spatial dispersion in a numerical decoder.}
A centered spatial stencil changes the free frequency
from $|k|$ to a known $\omega_h(k)$.
Even when $kh$ resolves the wave, a delayed slot incurs
phase error $d_k(\omega_h(k)-|k|)$.
For a stencil of order $p$ and all recorded modes
through $K$ with $Kh$ in the resolved range, a
sufficient condition to preserve the continuum phase
margin is $T K(Kh)^p\ll1$.
This differs from a requirement expressed only in
points per spatial wavelength.
For observations generated by a specified discretization
one can instead place the slots using its own
$\omega_h(k)$ in the two-phase rule.

A second option, when using a coarse inverse against
continuum-like data, corrects the modeled linear
transfer and retains the computed nonlinear remainder.
Write $\mathcal M_h(v)=C_h U(a)+D_h v+R_h(v)$,
where $C_h,D_h$ are the exact scalar damped-wave
transfers at the stencil frequencies, including the
observation weights. Replace the prediction by
\[
 \mathcal M_h^{\,\mathrm{corr}}(v)
    =\mathcal M_h(v)+(C-C_h)U(a)+(D-D_h)v .
\]
The correction uses the known initial displacement
and the current trial velocity.
If $\|R_h(v)-R_h(\widehat v)\|\le\epsilon_h\|v-\widehat v\|$
and the continuum diagonal $D$ has lower gain
$\gamma>\epsilon_h$, the correction makes the same
contraction proof apply to the inverse.
For a true continuum datum, the remaining error is
controlled by $R_h(v)-R(v)$ at the true initial
velocity, in addition to measurement noise.
Thus phase correction and nonlinear model bias are
separate. In the tests below the principal correction
removed the observed coarse-grid instability.
